\documentclass[a4paper, 11pt]{article}
\usepackage[utf8]{inputenc}
\usepackage[T1]{fontenc}
\usepackage[italian]{babel}
\usepackage{geometry}
\usepackage{graphicx}
\usepackage{tabularx}
\usepackage{fancyhdr}
\usepackage{hyperref}
\usepackage{enumitem} % Per gestire meglio la spaziatura degli elenchi
\usepackage{listings} % Utile per eventuali snippet di codice negli studi tecnologici
\usepackage{xcolor}
\usepackage{float}
\usepackage{titlesec}

% Margini
\geometry{top=3cm, bottom=3cm, left=2.5cm, right=2.5cm}

% ==============================================================================
% VARIABILI DEL DOCUMENTO (Modificare solo questi valori per ogni nuovo file)
% ==============================================================================
\newcommand{\DocTitle}{Piano di Progetto}
\newcommand{\DocSubtitle}{Capitolato C9 -- SpecCheck (Zucchetti S.p.A.)}
\newcommand{\DocVersion}{0.1.0}
\newcommand{\DocStatus}{Bozza} % Bozza | In verifica | Approvato
\newcommand{\DocDate}{2026/11/22}
\newcommand{\DocUsage}{Esterno} % Interno | Esterno
\newcommand{\DocRecipients}{0x06tudio \\ Zucchetti S.p.A. \\ Prof. Tullio Vardanega \\ Prof. Riccardo Cardin}
% ==============================================================================

% Comandi di utilità del gruppo
\newcommand{\issue}[1]{\href{https://github.com/0x06tudio/docs/issues/#1}{\textit{\##1}}}
\newcommand{\term}[1]{\textit{#1}$_G$} % Per eventuali termini da inserire nel Glossario

% Link ai verbali pubblicati nel repository (formato: verbale interno/esterno, AAAA, MM, GG)
\newcommand{\vint}[3]{\href{https://raw.githubusercontent.com/0x06tudio/docs/main/out/verbali/interni/vi_#1_#2_#3.pdf}{Verbale interno del #3/#2/#1}}
\newcommand{\vest}[3]{\href{https://raw.githubusercontent.com/0x06tudio/docs/main/out/verbali/esterni/ve_#1_#2_#3.pdf}{Verbale esterno del #3/#2/#1}}

% Link agli studi tecnologici pubblicati nel repository
\newcommand{\studio}[2]{\href{https://raw.githubusercontent.com/0x06tudio/docs/main/out/RTB/tecnologie_esplorative/#1}{#2}}

% Evidenzia i valori da verificare/sostituire prima dell'approvazione
\newcommand{\tbd}[1]{\colorbox{yellow!60}{#1}}

% Scostamenti rispetto al preventivo: rosso = peggio, verde = meglio
\newcommand{\dpos}[1]{\textcolor{red!80!black}{(+#1)}}
\newcommand{\dneg}[1]{\textcolor{green!50!black}{(-#1)}}
\newcommand{\rdec}[1]{\textcolor{red!80!black}{(-#1)}} % riduzione delle risorse rimanenti

% Paragrafi numerati su riga propria (4.2.1.1, 4.2.1.2, ...)
\setcounter{secnumdepth}{4}
\setcounter{tocdepth}{4}
\titleformat{\paragraph}[hang]{\normalfont\normalsize\bfseries}{\theparagraph}{1em}{}
\titlespacing*{\paragraph}{0pt}{1.5ex plus 1ex minus .2ex}{1ex}

% Configurazione stile link
\hypersetup{
    colorlinks=true,
    linkcolor=black,
    urlcolor=blue!70!black,
    pdftitle={\DocTitle},
    pdfauthor={0x06tudio}
}

% Configurazione stile blocchi di codice
\lstset{
    basicstyle=\ttfamily\small,
    breaklines=true,
    frame=single,
    backgroundcolor=\color{gray!5},
    rulecolor=\color{gray!40},
    captionpos=b
}

% Intestazione e piè di pagina
\pagestyle{fancy}
\fancyhf{}
\fancyhead[L]{\small \textit{0x06tudio} -- \DocTitle}
\fancyhead[R]{\small Anno Accademico 2026/2027}
\fancyfoot[L]{\small Versione \DocVersion}
\fancyfoot[R]{\small Pagina \thepage}
\renewcommand{\headrulewidth}{0.4pt}
\renewcommand{\footrulewidth}{0.4pt}

\begin{document}
\thispagestyle{plain}

% --- INTESTAZIONE PRIMA PAGINA ---
\begin{minipage}[t]{0.6\textwidth}
    \vspace{0pt}
    {\huge \textbf{0x06tudio}} \\[0.5em]
    {\large Progetto didattico di Ingegneria del Software} \\[0.3em]
    {\normalsize Università degli Studi di Padova}
\end{minipage}
\begin{minipage}[t]{0.4\textwidth}
    \raggedleft
    \vspace{0pt}
    \includegraphics[width=5.5cm]{../media/logo.png} \\[0.5em]
    \small
    \href{mailto:0x06tudio@gmail.com}{0x06tudio@gmail.com}
\end{minipage}

\vspace{1cm}

\begin{center}
    {\LARGE \textbf{\DocTitle}} \\[0.4em]
    {\large \textit{\DocSubtitle}}
\end{center}

\vspace{0.8cm}

% --- TABELLA INFORMAZIONI DEL DOCUMENTO ---
\section*{Informazioni sul documento}
\renewcommand{\arraystretch}{1.25}
\begin{tabularx}{\textwidth}{l X}
    \textbf{Versione:} & \DocVersion \\
    \textbf{Stato:} & \DocStatus \\
    \textbf{Data ultimo aggiornamento:} & \DocDate \\
    \textbf{Uso:} & \DocUsage \\
    \textbf{Redattori:} & Davide Colabove \\
                        & Catalin Chirilov \\
                        & Andrea Maggio \\
                        & Pietro Brocadello \\
                        & Alberto Cavalletto \\
                        & Nicolò Beltrame \\
    \textbf{Verificatori:} & Nicolò Beltrame \\
                           & Davide Colabove \\
                           & Catalin Chirilov \\
                           & Andrea Maggio \\
                           & Pietro Brocadello \\
                           & Alberto Cavalletto \\
    \textbf{Destinatari:} & \DocRecipients \\
\end{tabularx}

\vspace{0.5cm}

% --- REGISTRO DELLE MODIFICHE ---
\section*{Registro delle modifiche}
\renewcommand{\arraystretch}{1.25}
\begin{tabularx}{\textwidth}{|c|c|X|c|c|}
    \hline
    \textbf{Versione} & \textbf{Data} & \textbf{Descrizione} & \textbf{Autore} & \textbf{Verificatore} \\ \hline
    0.1.0 & 2026/11/22 & Stesura iniziale del documento (Sezioni 1, 2, 3, 4.1, 4.2.1 e inizio 4.2.2) & Pietro Brocadello & Nicolò Beltrame\\ \hline
\end{tabularx}

\newpage
\tableofcontents
\newpage
\listoftables
\newpage

% ==============================================================================
\section{Introduzione}

\subsection{Informazioni generali}
Il \term{Piano di Progetto} è il documento che descrive le attività svolte e da svolgere durante il \term{progetto} didattico di Ingegneria del Software relativo al capitolato C9 -- SpecCheck.

Viene utilizzato dal gruppo per la mitigazione dei rischi e per una corretta \term{pianificazione} dei ruoli e delle attività, con lo scopo di automiglioramento, seguendo un \term{processo} incrementale.

Le modalità operative e i dettagli relativi al \term{way of working} sono riportati nelle \term{Norme di Progetto}, consultabili nei Riferimenti normativi.

\subsection{Definizioni, acronimi e abbreviazioni}
Per tutti i termini tecnici, gli acronimi e le definizioni utilizzati nel documento si rimanda al \href{https://0x06tudio.github.io/docs/html/glossario.html}{Glossario}, disponibile come documento separato.

Ogni parola presente nel Glossario viene segnata come segue: \term{termine}.

\subsection{Riferimenti}

\subsubsection{Riferimenti normativi}
\begin{itemize}[noitemsep]
    \item \term{Capitolato} C9 -- SpecCheck -- Zucchetti S.p.A.: \\
    \url{https://www.math.unipd.it/~tullio/IS-1/2026/Progetto/C9-SpecCheck.pdf} \\
    Ultima consultazione: \tbd{22 novembre 2026};
    \item \term{Norme di Progetto} (versione \tbd{X.Y.Z}): \\
    \tbd{inserire il link al PDF nel repository}\\
    Ultima consultazione: \tbd{AAAA/MM/GG}.
\end{itemize}

\subsubsection{Riferimenti informativi}
\begin{itemize}[noitemsep]
    \item Preventivo dei costi e degli impegni orari (candidatura): \\
    \url{https://raw.githubusercontent.com/0x06tudio/docs/main/out/candidatura/preventivo_costi_e_impegni_orari.pdf} \\
    Ultima consultazione: 22 novembre 2026;
    \item Studio di fattibilità (candidatura): \\
    \url{https://raw.githubusercontent.com/0x06tudio/docs/main/out/candidatura/studio_di_fattibilita.pdf} \\
    Ultima consultazione: 22 novembre 2026;
    \item Tecnologie esplorative (cartella di approfondimenti tecnologici): \\
    \url{https://github.com/0x06tudio/docs/tree/main/out/RTB/tecnologie_esplorative} \\
    Ultima consultazione: 22 novembre 2026. In particolare:
    \begin{itemize}[noitemsep]
        \item Modulo 1 -- \studio{studio_tecnologico_harness_e_opencode.pdf}{Sistemi di Harness e Gestione del Contesto AI}, versione 0.1.0;
        \item Modulo 2 -- \studio{Benchmark_SpecCheck.pdf}{Benchmark per Agenti Software e Metodologie Spec-Driven}, versione 0.1.0;
        \item Modulo 3 -- \studio{ingestion_parsing.pdf}{Ingestion e Parsing di Diagrammi Architetturali}, versione 0.1.0;
        \item Modulo 4 -- \studio{tracciabilita_semantica.pdf}{Tracciabilità Semantica: Casi d'Uso vs Specifiche}, versione 0.1.0;
        \item Modulo 5 -- \studio{Studio_Tecnologico_Modulo5_AI_Code_Analysis.pdf}{Comprensione Comportamentale del Codice (AI Code Analysis)}, versione 0.1.0;
        \item Modulo 6 -- \studio{architettura_integrazione_tooling_dashboarding.pdf}{Architettura di Integrazione, Tooling e Dashboarding}, versione 0.1.0;
    \end{itemize}
    \item Glossario: \\
    \url{https://0x06tudio.github.io/docs/html/glossario.html} \\
    Ultima consultazione: 22 novembre 2026;
    \item Standard ISO/IEC 12207:1995: \\
    \url{https://www.math.unipd.it/~tullio/IS-1/2009/Approfondimenti/ISO_12207-1995.pdf} \\
    Ultima consultazione: 22 novembre 2026.
\end{itemize}

% ==============================================================================
\section{Analisi e gestione dei rischi}
In questa sezione il gruppo descrive l'analisi dei potenziali rischi associati al progetto.

Per mitigare tali rischi, abbiamo scelto di adottare un preventivo orario più contenuto all'inizio di ogni \term{sprint}, in particolare nelle fasi iniziali del progetto, quando la padronanza delle metodologie di lavoro è minore.

Ogni rischio è stato identificato ed analizzato, e sono state definite apposite strategie di mitigazione e controllo.

I rischi possono appartenere a tre macro-categorie:
\begin{itemize}[noitemsep]
    \item \textbf{RE (Rischi Economici):} rischi che influenzano tempi, costi, produttività o impegno orario del gruppo;
    \item \textbf{RP (Rischi Personali):} rischi legati ai singoli membri, alle loro attività, disponibilità, competenze e condizioni individuali;
    \item \textbf{RA (Rischi Ambientali):} rischi dovuti a fattori esterni, imprevisti o condizioni non controllabili direttamente dal gruppo.
\end{itemize}

Ciascun rischio è identificato da un codice che ne indica la categoria seguita da un numero incrementale (es.\ RE\#, RP\#, RA\#).

Ogni rischio possiede una frequenza probabile di avvenimento e un impatto, espressi su una scala a cinque livelli: \{Molto bassa, Bassa, Media, Elevata, Molto elevata\}.

\renewcommand{\arraystretch}{1.25}

\subsection{RE1: Rischi tecnologici}
\begin{table}[H]
\centering
\begin{tabularx}{\textwidth}{|l|X|}
    \hline
    \textbf{Tipologia} & \textbf{Descrizione} \\ \hline
    Codice & RE1 \\ \hline
    Nome & Tecnologie \\ \hline
    Descrizione & Il gruppo non ha esperienza con le tecnologie richieste dal capitolato: skill, comandi e plug-in per OpenCode, tool deterministici (Sidecar), dashboard di sintesi e integrazione con modelli linguistici. Sono inclusi possibili errori e lunghe sessioni di debugging. \\ \hline
    Mitigazione & All'inizio di ogni sprint si valuta il livello di conoscenza delle tecnologie, riducendo le ore produttive previste per includere il tempo di apprendimento. Le tecnologie vengono studiate in anticipo nei documenti di tecnologie esplorative e validate con un \term{Proof of Concept}. \\ \hline
    Frequenza probabile di avvenimento & Elevata \\ \hline
    Ripercussioni & Elevate \\ \hline
\end{tabularx}
\caption{Informazioni sul rischio RE1}
\end{table}

\subsection{RE2: Sovraccarico o stima errata del lavoro}
\begin{table}[H]
\centering
\begin{tabularx}{\textwidth}{|l|X|}
    \hline
    \textbf{Tipologia} & \textbf{Descrizione} \\ \hline
    Codice & RE2 \\ \hline
    Nome & Stima errata delle attività \\ \hline
    Descrizione & Le \term{task} possono richiedere più tempo del previsto, causando ritardi negli sprint e sovraccarico dei membri del gruppo. \\ \hline
    Mitigazione & Pianificare task realistiche e suddividere le task grandi in sotto-task. \\ \hline
    Frequenza probabile di avvenimento & Media \\ \hline
    Ripercussioni & Elevate \\ \hline
\end{tabularx}
\caption{Informazioni sul rischio RE2}
\end{table}

\subsection{RE3: Non determinismo dei modelli linguistici}
\begin{table}[H]
\centering
\begin{tabularx}{\textwidth}{|l|X|}
    \hline
    \textbf{Tipologia} & \textbf{Descrizione} \\ \hline
    Codice & RE3 \\ \hline
    Nome & Non determinismo e dipendenza dai modelli linguistici \\ \hline
    Descrizione & Il prodotto si appoggia a modelli linguistici il cui output può variare tra esecuzioni e la cui disponibilità o i cui costi d'uso possono cambiare, rendendo difficili test ripetibili e stime affidabili. \\ \hline
    Mitigazione & Separare la parte deterministica (tool del Sidecar) da quella inferenziale, vincolare gli output a schemi JSON, validarli con controlli automatici, verificare le skill anche sui modelli locali compatti (Gemma 4 26B, Qwen 3.8 27B) e documentare le scelte nel Piano di Qualifica. \\ \hline
    Frequenza probabile di avvenimento & Media \\ \hline
    Ripercussioni & Medie \\ \hline
\end{tabularx}
\caption{Informazioni sul rischio RE3}
\end{table}

\subsection{RP1: Attività universitarie}
\begin{table}[H]
\centering
\begin{tabularx}{\textwidth}{|l|X|}
    \hline
    \textbf{Tipologia} & \textbf{Descrizione} \\ \hline
    Codice & RP1 \\ \hline
    Nome & Attività universitarie \\ \hline
    Descrizione & Tutti i componenti del gruppo hanno almeno un altro esame oltre a questo corso di Ingegneria del Software. \\ \hline
    Mitigazione & All'inizio di ogni sprint si tiene conto dell'impegno richiesto dagli altri corsi, riducendo il preventivo orario per renderlo più realistico. \\ \hline
    Frequenza probabile di avvenimento & Media \\ \hline
    Ripercussioni & Elevate \\ \hline
\end{tabularx}
\caption{Informazioni sul rischio RP1}
\end{table}

\subsection{RP2: Attività extra-universitarie}
\begin{table}[H]
\centering
\begin{tabularx}{\textwidth}{|l|X|}
    \hline
    \textbf{Tipologia} & \textbf{Descrizione} \\ \hline
    Codice & RP2 \\ \hline
    Nome & Attività extra-universitarie \\ \hline
    Descrizione & Le attività extra-universitarie dei componenti (lavoro, sport, impegni personali) sono frequenti e riducono la disponibilità oraria. \\ \hline
    Mitigazione & All'inizio di ogni sprint si tiene conto degli impegni extra, riducendo il preventivo orario per renderlo coerente con la disponibilità reale. \\ \hline
    Frequenza probabile di avvenimento & Elevata \\ \hline
    Ripercussioni & Medie \\ \hline
\end{tabularx}
\caption{Informazioni sul rischio RP2}
\end{table}

\subsection{RP3: Incomprensioni}
\begin{table}[H]
\centering
\begin{tabularx}{\textwidth}{|l|X|}
    \hline
    \textbf{Tipologia} & \textbf{Descrizione} \\ \hline
    Codice & RP3 \\ \hline
    Nome & Incomprensioni \\ \hline
    Descrizione & Un membro può svolgere male una task per incomprensioni o per una definizione poco chiara del compito. \\ \hline
    Mitigazione & Il gruppo deve concordare task sempre chiare e ogni membro deve verificare di aver compreso correttamente cosa deve fare. \\ \hline
    Frequenza probabile di avvenimento & Molto bassa \\ \hline
    Ripercussioni & Elevate \\ \hline
\end{tabularx}
\caption{Informazioni sul rischio RP3}
\end{table}

\subsection{RA1: Imprevisti}
\begin{table}[H]
\centering
\begin{tabularx}{\textwidth}{|l|X|}
    \hline
    \textbf{Tipologia} & \textbf{Descrizione} \\ \hline
    Codice & RA1 \\ \hline
    Nome & Imprevisti \\ \hline
    Descrizione & È naturale che possano verificarsi imprevisti di tipo universitario, familiare o personale che influenzano la produttività. \\ \hline
    Mitigazione & La persona coinvolta deve avvisare tempestivamente il gruppo per permettere una ridistribuzione del lavoro. \\ \hline
    Frequenza probabile di avvenimento & Bassa \\ \hline
    Ripercussioni & Molto elevate \\ \hline
\end{tabularx}
\caption{Informazioni sul rischio RA1}
\end{table}

\subsection{RA2: Rischi legati ai requisiti}
\begin{table}[H]
\centering
\begin{tabularx}{\textwidth}{|l|X|}
    \hline
    \textbf{Tipologia} & \textbf{Descrizione} \\ \hline
    Codice & RA2 \\ \hline
    Nome & Requisiti instabili o incompleti \\ \hline
    Descrizione & I requisiti possono essere interpretati in modo errato o modificati durante il progetto, causando ritardi e lavoro non coerente con le aspettative. \\ \hline
    Mitigazione & Effettuare incontri chiarificatori con la \term{proponente} Zucchetti S.p.A., validare i requisiti dopo ogni sprint e produrre documentazione condivisa. \\ \hline
    Frequenza probabile di avvenimento & Media \\ \hline
    Ripercussioni & Elevate \\ \hline
\end{tabularx}
\caption{Informazioni sul rischio RA2}
\end{table}

% ==============================================================================
\section{Pianificazione nel lungo termine}

\subsection{Introduzione}
Come definito nel \href{https://raw.githubusercontent.com/0x06tudio/docs/main/out/candidatura/preventivo_costi_e_impegni_orari.pdf}{Preventivo dei costi e degli impegni orari}, il gruppo si impegna a concludere il progetto entro il 24/03/2027, con un costo complessivo massimo di 11.540\,\texteuro{} per un monte ore totale di 552 ore (circa 92 ore a componente).

Il preventivo iniziale, già definito nel documento menzionato precedentemente, è il seguente:

\begin{table}[H]
\centering
\begin{tabularx}{\textwidth}{|X|c|c|c|}
    \hline
    \textbf{Ruolo} & \textbf{Ore totali} & \textbf{Costo (\texteuro/h)} & \textbf{Costo totale (\texteuro)} \\ \hline
    \term{Responsabile} & 48 & 30 & 1440 \\ \hline
    \term{Amministratore} & 48 & 20 & 960 \\ \hline
    \term{Analista} & 105 & 25 & 2625 \\ \hline
    \term{Progettista} & 125 & 25 & 3125 \\ \hline
    \term{Verificatore} & 90 & 15 & 1350 \\ \hline
    \term{Programmatore} & 136 & 15 & 2040 \\ \hline
    \textbf{Totale} & \textbf{552} & -- & \textbf{11540} \\ \hline
\end{tabularx}
\caption{Ore per ciascun ruolo e calcolo dei costi totali}
\end{table}

La data di consegna prevista per la \term{RTB} è il 18/12/2026.

\subsection{Attività previste per la RTB}
Le attività previste per la \term{RTB} (\term{Requirements and Technology Baseline}) sono riassunte nella seguente tabella:

\begin{table}[H]
\centering
\begin{tabularx}{\textwidth}{|l|X|}
    \hline
    \textbf{Attività} & \textbf{Descrizione} \\ \hline
    \term{Analisi dei Requisiti} & Documento esterno che definisce i requisiti funzionali, non funzionali e di vincolo, e illustra i casi d'uso relativi a requisiti obbligatori, non obbligatori e di sistema. \\ \hline
    \term{Piano di Progetto} & Documento esterno che descrive le attività svolte e da svolgere durante il progetto. \\ \hline
    \term{Piano di Qualifica} & Documento esterno fondamentale per la gestione e il monitoraggio continuo della qualità del progetto e dei processi coinvolti nel suo ciclo di vita. \\ \hline
    \term{Norme di Progetto} & Documento interno che descrive nel dettaglio il way of working del gruppo. \\ \hline
    Glossario & Documento interno che facilita la comprensione dei termini tecnici nella comunicazione tra il gruppo e la proponente. \\ \hline
    \term{Proof of Concept} & Eseguibile che dimostra la coesione tra le tecnologie scelte. \\ \hline
    Tecnologie esplorative & Documenti di approfondimento sulle tecnologie del capitolato, a supporto delle scelte tecnologiche del Proof of Concept. \\ \hline
    Lettera di presentazione RTB & Simile alla Lettera di Presentazione per la Candidatura, esprime la volontà del gruppo di partecipare alla revisione per la RTB. \\ \hline
\end{tabularx}
\caption{Attività previste per la RTB}
\end{table}

Ogni attività che inizia, continua o finalizza la stesura delle documentazioni menzionate nella tabella è definita all'interno di ogni sprint nella sezione ``Attività da svolgere''.

\subsection{Attività previste per la PB}
Questa sezione verrà redatta dopo il superamento della \term{RTB}, in modo da consentire una stesura più accurata durante la \term{PB} (\term{Product Baseline}).

% ==============================================================================
\section{Pianificazione a breve termine}

\subsection{Introduzione}
Il gruppo 0x06tudio ha stabilito di procedere con un approccio \term{Agile} allo svolgimento del progetto, definendo la durata dello \term{sprint} di \tbd{2 settimane} per avere un'organizzazione migliore.

Seguendo questo principio, il gruppo si impegna a stabilire all'inizio di ogni sprint i ruoli all'interno del gruppo, che ruotano a ogni sprint, e le attività per le settimane successive, mantenendo la possibilità di modificarli qualora le esigenze organizzative lo imponessero. L'assegnazione dei ruoli di ciascuno sprint è consultabile nella pagina \href{https://0x06tudio.github.io/docs/html/gruppo.html}{Gruppo} del sito della documentazione.

Per verificare la continua e corretta prosecuzione delle attività, il gruppo rimane in contatto con la proponente Zucchetti S.p.A. tramite incontri telematici. Lo stato di avanzamento delle ore è consultabile nel \href{https://0x06tudio.github.io/docs/html/cruscotto.html}{Cruscotto di progetto}.

Seguiranno ora le descrizioni dei vari periodi di lavoro, nelle quali verranno esposte le sezioni:
\begin{enumerate}[noitemsep]
    \item Informazioni generali e attività da svolgere
    \item Rischi attesi
    \item Preventivo
    \item \term{Consuntivo}
    \item Aggiornamento delle risorse rimanenti
    \item Rischi incontrati
    \item \term{Efficacia} delle strategie di gestione di mitigazione dei rischi
    \item Migliorie da attuare per le attività future
    \item Retrospettiva
\end{enumerate}

\subsection{Requirements and Technology Baseline}

% ------------------------------------------------------------------------------
\subsubsection{Sprint 1}

\paragraph{Informazioni generali e attività da svolgere}
\begin{table}[H]
\centering
\begin{tabular}{|l|l|}
    \hline
    Inizio & 09/11/2026 \\ \hline
    Fine prevista & 22/11/2026 \\ \hline
    Fine reale & 22/11/2026 \\ \hline
    Giorni di ritardo & \tbd{0} \\ \hline
\end{tabular}
\end{table}

Questo primo periodo ha una durata di due settimane e comprende le attività di candidatura e di avvio del progetto.

Le attività da svolgere, definite nei verbali interni
\vint{2026}{10}{16}, \vint{2026}{10}{22}, \vint{2026}{10}{26},
\vint{2026}{10}{30}, \vint{2026}{11}{03}, \vint{2026}{11}{06} e
\vint{2026}{11}{13}, sono:
\begin{itemize}[noitemsep]
    \item Incontri con le proponenti e redazione dei verbali esterni (\vest{2026}{10}{20}, \vest{2026}{10}{22}, \vest{2026}{10}{26}, \vest{2026}{11}{13});
    \item Stesura del Resoconto dei capitolati e dello Studio di fattibilità;
    \item Stesura del Preventivo dei costi e degli impegni orari;
    \item Stesura della Lettera di presentazione;
    \item Creazione del repository della documentazione e del template LaTeX dei documenti;
    \item Creazione del sito della documentazione su GitHub Pages;
    \item Creazione del Glossario;
    \item Stesura dei diari di bordo 1--4;
    \item Inizio stesura delle Norme di Progetto.
\end{itemize}
\textit{Elenco delle attività verificato sui verbali interni ed esterni.}

\paragraph{Rischi attesi}
I possibili rischi attesi sono:
\begin{itemize}[noitemsep]
    \item RP1: alcuni componenti del gruppo sono impegnati con lo studio di altri corsi;
    \item RA1: un imprevisto a qualche componente;
    \item RE1: rischio dovuto all'impostazione di repository, pipeline e strumenti di documentazione;
    \item RE2: rischio di stima errata del lavoro;
    \item RP3: rischio di incomprensione del modo corretto con cui stendere i documenti;
    \item RA2: rischio di interpretazione errata del capitolato.
\end{itemize}

\paragraph{Preventivo}
Si prospetta l'utilizzo delle seguenti risorse:

\begin{table}[H]
\centering
\renewcommand{\arraystretch}{1.2}
\begin{tabular}{|l|c|c|c|c|c|c|}
    \hline
    & \rotatebox{60}{\term{Responsabile}} & \rotatebox{60}{\term{Amministratore}} & \rotatebox{60}{\term{Analista}} & \rotatebox{60}{\term{Programmatore}} & \rotatebox{60}{\term{Verificatore}} & \rotatebox{60}{\term{Progettista}} \\ \hline
    D. Colabove & \tbd{8} & -- & -- & -- & -- & -- \\ \hline
    C. Chirilov & -- & \tbd{10} & -- & -- & \tbd{1} & -- \\ \hline
    A. Maggio & -- & -- & \tbd{10} & -- & \tbd{1} & -- \\ \hline
    P. Brocadello & -- & -- & -- & -- & -- & \tbd{6} \\ \hline
    A. Cavalletto & -- & -- & -- & \tbd{4} & \tbd{1} & -- \\ \hline
    N. Beltrame & -- & -- & -- & -- & \tbd{8} & -- \\ \hline
\end{tabular}
\caption{Sprint 1: Preventivo}
\end{table}

Il totale preventivato per lo sprint è \tbd{\textbf{1065\,\texteuro}}.

\textit{\tbd{Ore illustrative: sostituire con le ore realmente preventivate nel verbale di inizio sprint.}}

\paragraph{Consuntivo}
Le risorse effettivamente utilizzate sono nella seguente tabella:

\begin{table}[H]
\centering
\renewcommand{\arraystretch}{1.2}
\begin{tabular}{|l|c|c|c|c|c|c|}
    \hline
    & \rotatebox{60}{\term{Responsabile}} & \rotatebox{60}{\term{Amministratore}} & \rotatebox{60}{\term{Analista}} & \rotatebox{60}{\term{Programmatore}} & \rotatebox{60}{\term{Verificatore}} & \rotatebox{60}{\term{Progettista}} \\ \hline
    D. Colabove & \tbd{9} \dpos{1} & -- & -- & -- & -- & -- \\ \hline
    C. Chirilov & -- & \tbd{8} \dneg{2} & -- & -- & \tbd{1} & -- \\ \hline
    A. Maggio & -- & -- & \tbd{12} \dpos{2} & -- & \tbd{1} & -- \\ \hline
    P. Brocadello & -- & -- & -- & -- & -- & \tbd{5} \dneg{1} \\ \hline
    A. Cavalletto & -- & -- & -- & \tbd{4} & \tbd{1} & -- \\ \hline
    N. Beltrame & -- & -- & -- & -- & \tbd{9} \dpos{1} & -- \\ \hline
\end{tabular}
\caption{Sprint 1: Consuntivo}
\end{table}

Il totale consuntivato per lo sprint è \tbd{\textbf{1095\,\texteuro}}.

\textit{\tbd{Ore illustrative: sostituire con le ore effettivamente rendicontate (vedi Cruscotto).}} Le ore di analisi sono state superiori al preventivo a causa del tempo dedicato a comprendere in profondità il capitolato e a confrontarsi con la proponente, mentre le ore di amministrazione sono risultate inferiori grazie alla riusabilità del template dei documenti.

\paragraph{Aggiornamento delle risorse rimanenti}

\begin{table}[H]
\centering
\renewcommand{\arraystretch}{1.2}
\resizebox{\textwidth}{!}{%
\begin{tabular}{|l|c|c|c|c|c|}
    \hline
    \textbf{Ruolo} & \textbf{Costo} & \textbf{Ore} & \textbf{Costo effettivo} & \textbf{Ore rimanenti} & \textbf{Budget rimanente} \\ \hline
    \term{Responsabile} & 30 \texteuro/h & 9 & 270 & 39 \rdec{9} & 1170 \rdec{270} \\ \hline
    \term{Amministratore} & 20 \texteuro/h & 8 & 160 & 40 \rdec{8} & 800 \rdec{160} \\ \hline
    \term{Analista} & 25 \texteuro/h & 12 & 300 & 93 \rdec{12} & 2325 \rdec{300} \\ \hline
    \term{Progettista} & 25 \texteuro/h & 5 & 125 & 120 \rdec{5} & 3000 \rdec{125} \\ \hline
    \term{Verificatore} & 15 \texteuro/h & 12 & 180 & 78 \rdec{12} & 1170 \rdec{180} \\ \hline
    \term{Programmatore} & 15 \texteuro/h & 4 & 60 & 132 \rdec{4} & 1980 \rdec{60} \\ \hline
    \textbf{Totale} & -- & 50 & 1095 & 502 \rdec{50} & 10445 \rdec{1095} \\ \hline
\end{tabular}}
\caption{Sprint 1: Aggiornamento risorse}
\end{table}

\paragraph{Rischi incontrati}
Sono stati incontrati i rischi \tbd{RE1 ed RE2}, già definiti nei rischi attesi. Il primo si è manifestato nella configurazione di repository, pipeline di compilazione dei documenti e sito della documentazione, più laboriosa del previsto. Il secondo si è manifestato nelle ore di analisi, sottostimate in fase di preventivo.

\paragraph{Efficacia delle strategie di gestione di mitigazione dei rischi}
Per il rischio RE1, trattandosi della prima esperienza del gruppo con un progetto di questa complessità, si è scelto di riservare le prime attività alla predisposizione degli strumenti di lavoro, in modo da ridurre i problemi nelle fasi successive.

Per il rischio RE2, il gruppo cerca di migliorare sprint dopo sprint la stesura del preventivo, per renderlo più accurato.

\paragraph{Migliorie da attuare per le attività future}
Nei prossimi sprint il gruppo si pone l'obiettivo di proseguire la stesura dei documenti rimanenti (Analisi dei Requisiti, Norme di Progetto, Piano di Progetto e Piano di Qualifica) e di avviare lo studio delle tecnologie del capitolato.

\paragraph{Retrospettiva}
Al termine dello sprint \tbd{sono state completate tutte le attività inserite tra quelle da svolgere}.

% ------------------------------------------------------------------------------
\subsubsection{Sprint 2}

\paragraph{Informazioni generali e attività da svolgere}
\begin{table}[H]
\centering
\begin{tabular}{|l|l|}
    \hline
    Inizio & 23/11/2026 \\ \hline
    Fine prevista & 06/12/2026 \\ \hline
    Fine reale & -- \\ \hline
    Giorni di ritardo & -- \\ \hline
\end{tabular}
\end{table}

Le attività da svolgere, definite nel \vint{2026}{11}{13} e nel \vest{2026}{11}{13}, sono:
\begin{itemize}[noitemsep]
    \item Inizio stesura Piano di Progetto;
    \item Continuazione stesura Norme di Progetto;
    \item Inizio individuazione Casi d'Uso e Requisiti;
    \item Stesura degli studi tecnologici dei Moduli 1--6: \studio{studio_tecnologico_harness_e_opencode.pdf}{Harness e Gestione del Contesto}, \studio{Benchmark_SpecCheck.pdf}{Benchmark e Metodologie Spec-Driven}, \studio{ingestion_parsing.pdf}{Ingestion e Parsing}, \studio{tracciabilita_semantica.pdf}{Tracciabilità Semantica}, \studio{Studio_Tecnologico_Modulo5_AI_Code_Analysis.pdf}{AI Code Analysis}, \studio{architettura_integrazione_tooling_dashboarding.pdf}{Integrazione, Tooling e Dashboarding};
    \item Aggiornamento Glossario;
    \item Stesura diari di bordo \tbd{5 e 6}.
\end{itemize}
\textit{\tbd{Verificare e correggere l'elenco delle attività in base ai verbali.}}

\paragraph{Rischi attesi}
I possibili rischi attesi sono:
\begin{itemize}[noitemsep]
    \item RE1: rischio dovuto alla poca conoscenza delle tecnologie da affrontare (OpenCode, Sidecar, modelli linguistici);
    \item RE2: sovraccarico o stima errata del lavoro;
    \item RP1: impegno dovuto allo studio di altri corsi;
    \item RP3: rischio di incomprensione di qualche task;
    \item RA2: rischio di interpretazione errata dei requisiti del capitolato.
\end{itemize}

\paragraph{Preventivo}
Si prospetta l'utilizzo delle seguenti risorse:

\begin{table}[H]
\centering
\renewcommand{\arraystretch}{1.2}
\begin{tabular}{|l|c|c|c|c|c|c|}
    \hline
    & \rotatebox{60}{\term{Responsabile}} & \rotatebox{60}{\term{Amministratore}} & \rotatebox{60}{\term{Analista}} & \rotatebox{60}{\term{Programmatore}} & \rotatebox{60}{\term{Verificatore}} & \rotatebox{60}{\term{Progettista}} \\ \hline
    D. Colabove & -- & \tbd{5} & -- & -- & \tbd{1} & -- \\ \hline
    C. Chirilov & \tbd{4} & -- & -- & -- & \tbd{1} & -- \\ \hline
    A. Maggio & -- & -- & \tbd{6} & -- & \tbd{1} & -- \\ \hline
    P. Brocadello & -- & -- & -- & \tbd{4} & \tbd{1} & -- \\ \hline
    A. Cavalletto & -- & -- & -- & \tbd{4} & \tbd{1} & -- \\ \hline
    N. Beltrame & -- & -- & \tbd{2} & -- & \tbd{2} & -- \\ \hline
\end{tabular}
\caption{Sprint 2: Preventivo}
\end{table}

Il totale preventivato per lo sprint è \textbf{650\,\texteuro}.

\textit{Assegnazione dei ruoli dello sprint 2.}

\end{document}
